%------------------------------------------------------------------------------%
\begin{question}
  Verifique se a matriz
  \[
    B=
    \begin{bmatrix}
      2 & 1 & 0 \\
      0 & 2 & 0 \\
      0 & 0 & 5
    \end{bmatrix}
  \]
  é diagonalizável
\end{question}

%------------------------------------------------------------------------------%
\begin{resolution}{Autovalores}
  Como a matriz é triangular
  \[
    p_{B}(\lambda)=(2-\lambda)^2(5-\lambda)
  \]
  Os autovalores são
  \[
    \lambda_{1}=2
    \qquad
    \lambda_{2}=5
  \]
  As multiplicidades algébricas são
  \[
    m_{a}(2)=2
    \qquad
    m_{a}(5)=1
  \]
\end{resolution}

%------------------------------------------------------------------------------%
\begin{resolution}{Autoespaços}
  Para \(\lambda=2\)
  \[
    (B-2I)\vec{v}=\vec{0}
  \]
  implica
  \[
    y=0
    \qquad
    z=0
  \]
  Logo
  \[
    E_{2}=\operatorname{span}
    \left\{
    \begin{bmatrix}
    1\\0\\0
    \end{bmatrix}
    \right\}
  \]
  Assim
  \[
    \dim(E_{2})=1
  \]
\end{resolution}

%------------------------------------------------------------------------------%
\begin{resolution}{Conclusão}
  Para \(\lambda=5\)
  \[
    E_{5}=\operatorname{span}
    \left\{
    \begin{bmatrix}
      0 \\
      0 \\
      1
    \end{bmatrix}
    \right\}
  \]
\end{resolution}

%------------------------------------------------------------------------------%
\begin{resolution}{Conclusão}
  Portanto, a soma das dimensões dos autoespaços é
  \[
    \dim(E_{2})+\dim(E_{5})=1+1=2
  \]
  Para diagonalizar uma matriz de ordem \(3\),

  precisamos de três autovetores linearmente independentes.

  Logo \(B\) não é diagonalizável
\end{resolution}

%------------------------------------------------------------------------------%
\endinput

